\scanpage{185}
Ceci suggère \deleted{d'associer à $\dot x$} de définir $\unsure{\varphi|\psi}$ par
\[
 \boxed{\psi(\dot x\wedge\dot x^2\wedge\text{--}\wedge\dot x^{p-1})
 =d\widetilde x^{\otimes\frac{p-1}{2}}\otimes 1_\Lambda}
 \tag*{($\ast\ast\ast$)}
\]
Ce qui donne une racine carrée de l'isom $(\ast\ast)$, et dont \unsure{revient|restent|restait} à prouver que $\psi$ est indépendant du choix de la $p$-base $\{x\}$. \deleted{Soit $y$ une autre $p$-base,}
\[
 \deleted{y=a_{p-1}x^{p-1}\mathrel{\text{---}}+a_0\qquad(a_i\in\Lambda)}
\]
\deleted{d'où}
\[
 \deleted{dy=[(p-1)a_{p-1}x^{p-2}\mathrel{\text{---}}+a_1]\,dx,}
\]
\deleted{donc on a}
\[
 \deleted{(p-1)a_{p-1}x^{\unsure{p-1|p-2}}\mathrel{\text{---}}+a_1\in{\unsure{\mathbb B|B}}^\ast}
\]
Or, le passage de l'un à l'autre se fait par un signe. On trouve un \deleted{\unsure{hom}} \inserted{morphisme} du schéma \inserted{$P$} [sur \deleted{\unsure{$\Lambda$}} \inserted{$Z=(X^{p/S},\Lambda)$}] des $p$-bases relatives, dans le revêtemt trivial à deux feuillets $(\mathbb Z/2\mathbb Z)_Z$ de $Z$. \deleted{\unsure{Cette ap}} Or \deleted{\unsure{$Z$}} $P$ sur $Z$ est à fibres connexes, et on voit donc que \deleted{cette application est} ce morphisme est \unsure{const\textsuperscript{ant}|const.|constante} fibre à fibre, ce qui donne ce qu'on veut $\ldots$.

Application : la dualité en cohomologie globale ($X$ lisse propre sur $S$ de dim rel. 1.)

\margin{?} Généralisation des relations précédentes à la dim supérieure ?
% archive sheet number: 184
